\documentclass[11pt]{article}
\usepackage[margin=1in]{geometry}
\usepackage{amsmath,amssymb,amsthm,microtype,booktabs}
\usepackage[hidelinks]{hyperref}
\allowdisplaybreaks[2]

\newtheorem{theorem}{Theorem}[section]
\newtheorem{lemma}[theorem]{Lemma}
\newcommand{\e}{\mathrm e}
\newcommand{\ii}{\mathrm i}
\newcommand{\dd}{\mathrm d}
\newcommand{\R}{\mathbb R}
\newcommand{\bits}{\{0,1\}^L}
\newcommand{\boxD}{\mathcal D}

\title{A quantics tensor train for the Newman heat flow\\
       \large Rank and observation error on a fixed complex box}
\author{Edward Baker}
\date{14 November 2026}

\begin{document}
\maketitle

\begin{abstract}
We bound the tensor-train ranks needed to approximate a finite source
state for the de Bruijn--Newman heat flow. The construction uses the
complete theta source, a Chebyshev approximation at the initial time,
and a Taylor polynomial for the heat multiplier. On
\(0\le\lambda\le1/2\), \(40\le\operatorname{Re} z\le80\),
\(\lvert\operatorname{Im} z\rvert\le8\), a 72-bit midpoint
discretization has a polynomial replacement with bond dimensions at most
\(641\). Its exact cosine contraction approximates the continuum
field in an explicitly scaled absolute norm with error below \(2^{-36}\).
The scalar integral of the same polynomial satisfies the same bound
with the midpoint contribution removed. The construction uses exact
coefficients and contractions.
\end{abstract}

\section{Introduction}

In the source coordinate \(u\), heat evolution is multiplication by a
Gaussian. Observation of the resulting field involves an oscillatory
integral, in which a positive source can undergo substantial cancellation.
A useful approximation statement must therefore control the scalar
observation as well as the source representation. We use the reversed
time \(\tau=1/2-\lambda\) and the scaled absolute error
\[
 \|E\|_{S,\boxD}=\sup_{(\lambda,z)\in\boxD}S(x)|E_\lambda(z)|,
 \qquad
 S(x)=\frac{\exp(\pi x/8)}{(1+x)^{7/4}},
 \qquad x=\operatorname{Re}z,
\]
where \(\boxD\) is the box in the abstract. Stirling's formula
\cite[Eq.~5.11.9]{dlmf} gives the scale \(x^{7/4}\exp(-\pi x/8)\)
for the prefactor of \(\zeta\) in the time-zero real-axis formula
derived below, motivating the reciprocal scale \(S(x)\). The error is
defined at zeros and uses the same real scale at each time and offset.

Binary indexing permits a sampled polynomial to be evaluated by a
chain of small matrices. Lindsey's interpolation construction gives
polynomial exactness for quantics tensor trains \cite[\S4.1]{lindsey}.
Here we use a direct polynomial translation formula and derive explicit
constants for the complete theta source. The tensor format represents
the midpoint state, and comparison with the exact integral of the same
polynomial isolates the cost in error of passing to that state.
The proofs below require source analyticity on a specified ellipse,
uniform bounds for the complete source and heat multiplier, and the
literal cosine observation.

\section{The complete source and analytic bounds}
\label{sec:source}

Write \(r_n(v)=\pi n^2\e^{4v}\). The source and its heat weight are
\begin{equation}
 \Phi(u)=\e^u\sum_{n\ge1}(2r_n(u)^2-3r_n(u))\e^{-r_n(u)},
 \qquad W_\lambda(u)=\e^{\lambda u^2}\Phi(u).
 \label{eq:source}
\end{equation}
For the reflection, let
\(\vartheta(t)=\sum_{n\in\mathbb Z}\e^{-\pi n^2t}\) for \(t>0\) and
\(p(u)=\e^u(\vartheta(\e^{4u})-1)/2\). The Jacobi transform
\(\vartheta(t)=t^{-1/2}\vartheta(1/t)\) follows by setting the theta
argument to zero and the lattice parameter to \(\ii t\) in
\cite[Eqs.~20.2.3 and 20.7.32]{dlmf}. It gives
\(p(-u)-p(u)=\sinh u\), so \(p'(0)=-1/2\). Differentiating twice,
term by term on compact real intervals, gives \((p''-p)/8=\Phi\).
The relation for \(p\) then implies evenness of the complete source.
Since each summand in \eqref{eq:source} is positive on \(u\ge0\),
\(\Phi\) is positive on the full line. The field is
\begin{equation}
 H_\lambda(z)=\frac12\int_{\R}W_\lambda(u)\e^{\ii zu}\,\dd u
             =\int_0^\infty W_\lambda(u)\cos(zu)\,\dd u.
 \label{eq:field}
\end{equation}
The weights in \eqref{eq:field} retain the total mass of the source.

For the normalization at \(\lambda=0\), put \(s=(1-\ii z)/2\) and
\(J(z)=\int_0^\infty p(u)\cos(zu)\,\dd u\).
For \(\operatorname{Re}s>1\), termwise Mellin integration gives
\[
 \pi^{-s/2}\Gamma(s/2)\zeta(s)
  =\int_0^\infty \frac{\vartheta(t)-1}{2}\,t^{s/2-1}\,\dd t.
\]
Splitting at \(1\) and using theta reflection below \(1\) gives the
sum of \(1/[s(s-1)]\) and
\[
 \int_1^\infty \frac{\vartheta(t)-1}{2}
       \bigl(t^{s/2-1}+t^{(1-s)/2-1}\bigr)\,\dd t=8J(z).
\]
The last identity uses \(t=\e^{4u}\). Multiplying by \(s(s-1)/2\)
continues the completed function \(\xi(s)\) to all \(s\). Two integrations
by parts, with \(p'(0)=-1/2\), give \(8H_0(z)=1/2-(1+z^2)J(z)=\xi(s)\).
All integrals over the positive \(u\) tail in this calculation converge
by the double exponential from \(r_n\).

Let \(R\) be the complex rectangle with
\[
  -1/16\le\operatorname{Re}v\le33/16,\qquad
  |\operatorname{Im}v|\le9/40 .
\]
We use this rectangle both for approximation on \([0,2]\) and for
estimates on the second derivative in midpoint quadrature.

\begin{lemma}[Complex source and integrand bounds]
\label{lem:complex}
The series for \(W_\lambda\) is holomorphic on a neighborhood of \(R\)
and \(\sup_R |W_\lambda|\le 2^{15}\) for \(0\le\lambda\le1/2\). On
\(\boxD\times[0,2]\),
\[
 \left|\frac{\dd^2}{\dd u^2}\bigl(W_\lambda(u)\cos(zu)\bigr)\right|
 \le M_2,\qquad M_2=2^{57}.
\]
\end{lemma}
\begin{proof}
For \(v\in R\) set \(t=|r_n(v)|\). Then
\(t/n^2\ge\pi\e^{-1/4}>2\) and
\(\operatorname{Re}r_n(v)=t\cos(4\operatorname{Im}v)>t/2\).
The second bound uses \(\cos a\ge1-a^2/2\) for \(|a|\le9/10\).
Maximizing \(t^2\e^{-t/4}\) and \(t\e^{-t/4}\) on \(t\ge0\) gives
\[
 |(2r_n(v)^2-3r_n(v))\e^{-r_n(v)}|
 \le (2t^2+3t)\e^{-t/2}
 \le (128/\e^2+12/\e)\e^{-t/4}<38\e^{-n^2/2}.
\]
The strict margins in the bounds for \(t/n^2\) and
\(\operatorname{Re}r_n/t\) persist on an enlarged rectangle,
so the series converges normally there. Also
\(\sum_{n\ge1}\e^{-n^2/2}\le 1/(\e^{1/2}-1)\le2\).
On \(R\), the prefactor in \eqref{eq:source} has
modulus at most \(\exp(33/16+(33/16)^2/2)<\e^5<2^8\).
Here and below \(\e^2<8\) follows from \(\e<11/4\), obtained by
bounding the exponential series tail from \(1/3!\) by a geometric series
of ratio \(1/4\). Thus \(76\cdot2^8<2^{15}\) bounds the complete
source on \(R\).

For \(u\in[0,2]\) the circle \(|v-u|=1/16\) lies in \(R\), and
\(|\operatorname{Im}(zv)|\le80/16+8(33/16)<22\).
Using \(|\cos(zv)|\le \exp(|\operatorname{Im}(zv)|)<2^{33}\),
Cauchy's derivative formula gives \(M_2=2!\,16^2\cdot2^{15+33}\).
\end{proof}

\begin{lemma}[Continuum tail and scale]
\label{lem:tail}
With \(S_*=\max_{x\in[40,80]}S(x)\) and
\(T=\int_2^\infty \e^{8u}W_{1/2}(u)\,\dd u\), we have
\[
 T<2^{-192},\qquad S_*\le2^{48}.
\]
In particular the tail of \eqref{eq:field} beyond \(2\) is at most \(T\)
in absolute value on \(\boxD\).
\end{lemma}
\begin{proof}
Set \(r=\pi\e^{4u}\) for real \(u\ge0\). The inequalities
\(\log n\le n-1\) and \(\pi(n^2-1)\ge6(n-1)\) give
\(n^4\e^{-\pi(n^2-1)}\le \e^{-2(n-1)}\).
Since \(r\ge\pi\), summing \eqref{eq:source} with
\(2r_n^2-3r_n\le2n^4r^2\) then gives
\[
 W_{1/2}(u)\le4\e^{u^2/2+u}r^2\e^{-r}.
\]
For \(u\ge2\),
\[
 u^2/2+9u+2\log r
 <u^2/2+17u+4
 \le \tfrac32\e^{4u}<r/2.
\]
For the middle inequality,
\(\tfrac32\e^{4u}\ge\tfrac32+6u+12u^2+16u^3\), whose cubic
lower bound exceeds \(u^2/2+17u+4\) for \(u\ge2\).
Thus the substitution \(s=r\), \(\dd u=\dd s/(4s)\), gives
\[
 T\le\int_{\pi\e^8}^{\infty}\e^{-s/2}\frac{\dd s}{s}
 \le\frac{2}{\pi\e^8}\e^{-\pi\e^8/2}<2^{-192}.
\]
For the last step \(\pi\e^8/2>(3/2)2^8>192\).
Because \(W_\lambda\le W_{1/2}\) and
\(|\cos(zu)|\le\e^{8u}\) on \(\boxD\), this bound applies to the
complex observation. Finally \(S_*\le\e^{10\pi}<\e^{32}<2^{48}\).
\end{proof}

The real-axis bound in the last proof, with the Gaussian factor removed,
dominates \(\Phi\) by a constant times
\(\e^{|u|}(\pi\e^{4|u|})^2\exp(-\pi\e^{4|u|})\).
Even after multiplication by a fixed power of \(|u|\) and by
\(\e^{A u^2+B|u|}\) for fixed \(A,B\ge0\), the result is integrable.
This gives local uniform majorants for \eqref{eq:field} and its partial
derivatives in \(\lambda,z\), so
\(\partial_\lambda H_\lambda=-\partial_z^2 H_\lambda\).
In time \(\tau=1/2-\lambda\), the observation obeys the forward heat
equation and the source satisfies \(W_\lambda(u)=\e^{-\tau u^2}W_{1/2}(u)\).

\section{A polynomial approximation to the heat orbit}
\label{sec:polynomial}

Let \(T_k\) be the Chebyshev polynomial determined by
\(T_k(\cos\theta)=\cos(k\theta)\).
Take \(\rho=5/4\) and \(G(w)=W_{1/2}(1+(w+w^{-1})/2)\).
For \(|w|=\rho\), the image ellipse has horizontal extent
\([-1/40,81/40]\) and imaginary semiaxis \(9/40\). The images
of all circles \(\rho^{-1}\le |w|\le\rho\) lie inside this
closed ellipse. By Lemma~\ref{lem:complex}, \(G\) has a Laurent series
on an open annulus containing these circles. If \(b_k\) are its
coefficients, symmetry \(G(w)=G(w^{-1})\) gives \(b_{-k}=b_k\);
Cauchy's formula on \(|w|=\rho\) gives
\(|b_k|\le M_\rho\rho^{-k}\) for \(k\ge0\), with \(M_\rho=2^{15}\).
Hence on \([0,2]\) the series is
\(b_0+2\sum_{k\ge1}b_kT_k(u-1)\). Define \(p_d\) by
truncating that sum at \(k=d\). The coefficients refer to the
complete source in \eqref{eq:source}, and the geometric tail gives
\[
 \|W_{1/2}-p_d\|_{\infty,[0,2]}\le
 \epsilon_d=\frac{2M_\rho\rho^{-d}}{\rho-1}.
\]

For \(0\le\tau\le1/2\), let
\[
 q_m(\tau,u)=\sum_{j=0}^m\frac{(-\tau u^2)^j}{j!},
 \qquad P_\tau(u)=p_d(u)q_m(\tau,u).
\]
Taylor's theorem for \(a\mapsto\e^{-a}\), whose derivatives have modulus
at most one on \(0\le a\le2\), gives
\(\|\e^{-\tau u^2}-q_m(\tau,u)\|_{\infty,[0,2]}\le E_m=2^{m+1}/(m+1)!\).
Multiplying the two approximations and writing their error as
\((W_{1/2}-p_d)\e^{-\tau u^2}+p_d(\e^{-\tau u^2}-q_m)\) yields
\begin{equation}
 \|W_\lambda-P_\tau\|_{\infty,[0,2]}\le
 \delta=\epsilon_d+(M_0+\epsilon_d)E_m,\qquad M_0=2^{15}.
 \label{eq:delta}
\end{equation}
Here \(M_0\) bounds \(W_{1/2}\) on \([0,2]\); the degree of \(P_\tau\)
is at most \(D_p=d+2m\).

\section{Binary states, the heat operator, and the readout}
\label{sec:tensors}

Take \(2^L\) cells on \([0,2]\) with width \(h=2/2^L\). Use bits
\(\boldsymbol b=(b_1,\dots,b_L)\in\bits\), in order from coarse
to fine. The midpoint is \(u_{\boldsymbol b}=a+\sum_{k=1}^L c_k b_k\),
where \(a=h/2\) and \(c_k=h2^{L-k}\).
The sampled state has amplitudes \(\sqrt h\,W_\lambda(u_{\boldsymbol b})\).
Since \(h\,2^L=2\), replacing these amplitudes by
\(\sqrt h\,P_\tau(u_{\boldsymbol b})\) gives a state error of at most
\(\sqrt{2}\,\delta\) in Euclidean norm by \eqref{eq:delta}.
A quantics tensor train (QTT) for any such array is a factorization
whose \(\boldsymbol b\)-entry is a scalar product of matrices
\(A_1(b_1)\cdots A_L(b_L)\). The first row and last column dimensions
are one; the intermediate dimensions bound the ranks across the
corresponding bit cuts.

\subsection{Translation gives exact polynomial cores}

Let \(v(u)=(1,u,\ldots,u^{D_p})^{\mathsf T}\).
For \(0\le j,k\le D_p\), set \(B(t)_{jk}=\binom jk t^{j-k}\)
if \(k\le j\) and \(0\) otherwise, with diagonal entries one.
The binomial theorem gives \(B(t)v(u)=v(u+t)\). If
\(\boldsymbol p_\tau\) is the row of monomial coefficients of \(P_\tau\),
successive translations give
\[
 \sqrt h\,P_\tau(u_{\boldsymbol b})
   =\sqrt h\,\boldsymbol p_\tau
     B(c_Lb_L)\cdots B(c_1b_1)v(a).
\]
Transposing the scalar puts the cores \(B(c_kb_k)^{\mathsf T}\)
in coarse-to-fine order. The initial row \(\sqrt h\,v(a)^{\mathsf T}\)
and final column \(\boldsymbol p_\tau^{\mathsf T}\) can be folded into
the first and last cores. Thus every internal bond has dimension
at most \(D_p+1\), irrespective of \(L\).

\subsection{Generator and finite-time multiplier}

On the finite bit space let \(N_k\) be multiplication by \(b_k\) at
site \(k\), \(X=aI+\sum_k c_kN_k\), and \(K=X^2\).
The exact state evolves by \(\partial_\tau\psi=-K\psi\).
The generator has a representation as a product of operator-valued
cores of bond dimension at most three,
\[
 K=(1,2a,a^2)\prod_{k=1}^{L}
      \begin{pmatrix}
       I_k&2c_kN_k&c_k^2N_k\\
       0&I_k&c_kN_k\\
       0&0&I_k
      \end{pmatrix}
       \begin{pmatrix}0\\0\\1\end{pmatrix}.
\]
Here products of entries on different sites are tensor products and
\(I_k\) denotes the local identity. On the bit basis the running row
\((1,2s,s^2)\) updates to \((1,2(s+c_k b_k),(s+c_k b_k)^2)\).
This proves the representation. Equivalently,
\[
 K=a^2I+\sum_k(2ac_k+c_k^2)N_k
    +2\sum_{j<k}c_jc_kN_jN_k.
\]
The terms commute, so \(\e^{-\tau K}\) factors exactly into their
exponentials. In the polynomial construction \(q_m(\tau,X)\) acts
diagonally on the state sampled from \(p_d\), and
\(\|\e^{-\tau K}-q_m(\tau,X)\|_{2\to2}\le E_m\) by the scalar
heat bound.

There is a different rank behavior for the exact heat multiplier
when \(\tau>0\). At a cut after site \(k\), write the midpoint index as
\(pM+q\), with \(M=2^{L-k}\), \(0\le p<2^k\), \(0\le q<M\), and
write \(\alpha_p=a+hMp\). The array of multiplier values has entries
\[
 \e^{-\tau(\alpha_p+hq)^2}
 =\e^{-\tau\alpha_p^2}\e^{-\tau h^2q^2}
    \bigl(\e^{-2\tau h\alpha_p}\bigr)^q .
\]
Removing the nonzero row and column factors gives a Vandermonde
matrix with distinct nodes. The first \(r=\min(2^k,M)\) rows and
columns have nonzero determinant, so the array has rank \(r\) across
the cut. This rank is for the exact multiplier in the stated bit
ordering. The error bound for the sampled state instead uses the
polynomial degree and the uniform Taylor remainder.

\subsection{The cosine contraction}

Each exponential \(\e^{\pm\ii z u_{\boldsymbol b}}\) is the product
\(\e^{\pm\ii za}\prod_k\e^{\pm\ii z c_k b_k}\). For a train with cores
\(A_k(b_k)\), the two contractions are
\begin{equation}
 \sqrt h\,\e^{\pm\ii za}
     \prod_{k=1}^{L}\bigl(A_k(0)+\e^{\pm\ii z c_k}A_k(1)\bigr).
 \label{eq:readout}
\end{equation}
For the polynomial state their half-sum is
\begin{equation}
 Q_{\lambda,L}(z)
 =h\sum_{\boldsymbol b\in\bits}P_\tau(u_{\boldsymbol b})
                               \cos(zu_{\boldsymbol b}).
 \label{eq:Q}
\end{equation}
The first \(\sqrt h\) belongs to the state and the second to the
observation. The half-sum implements the cosine on the positive
half-line, where the full-line Fourier factor \(1/2\) is already
accounted for by reflection. Formula \eqref{eq:readout} defines the
contraction without forming the dense array.

\section{Fixed-box rank and observation error}
\label{sec:theorem}

\begin{theorem}
Take \(d=512\), \(m=64\), and \(L=72\) in the construction above.
The polynomial midpoint state has QTT ranks at most \(641\) and
\[
 \|H-Q_{\cdot,L}\|_{S,\boxD}<2^{-36}.
\]
The scalar polynomial integral
\(I_\lambda(z)=\int_0^2P_\tau(u)\cos(zu)\,\dd u\)
satisfies the same bound with the midpoint contribution deleted.
\end{theorem}
\begin{proof}
The rank bound is \(512+2\cdot64+1\) by the translation construction.
Let \(f(u)=W_\lambda(u)\cos(zu)\). The integrated Taylor remainder
on a cell centered at \(u_{\boldsymbol b}\) has absolute value
at most \(M_2h^3/24\): the linear term integrates to zero, and the
pointwise remainder is bounded by
\(M_2|u-u_{\boldsymbol b}|^2/2\).
Summing over \(2/h\) cells bounds the quadrature error for the
complete integrand by \(2h^2M_2/24\). At the nodes,
\eqref{eq:delta} and \(|\cos(zu_{\boldsymbol b})|\le\e^{16}\)
bound the error from replacement by the polynomial by \(2\e^{16}\delta\).
Adding the tail of Lemma~\ref{lem:tail} then gives
\begin{equation}
 \|H-Q_{\cdot,L}\|_{S,\boxD}
 \le S_*\left(T+\frac{2h^2}{24}M_2+2\e^{16}\delta\right).
 \label{eq:transfer}
\end{equation}

The inequalities \((4/5)^4<1/2\) and
\(65!\ge2^{270}\) imply \(\epsilon_d<2^{-110}\),
\(E_m\le2^{-205}\). For the factorial bound, replace each factor by
the greatest power of two below or equal to it. The blocks
\(2\)--\(3\), \(4\)--\(7\), \(8\)--\(15\), \(16\)--\(31\),
\(32\)--\(63\), and \(64\)--\(65\) contribute exponents
\(2,8,24,64,160,12\), respectively. Also
\(2\e^{16}<2^{25}\), \(M_0+\epsilon_d<2^{16}\) and \(h=2^{-71}\).
With \(S_*\le 2^{48}\), the contributions in \eqref{eq:transfer} obey
\[
 \begin{array}{lr}
 \toprule
 \text{Contribution after scaling} & \text{Upper bound}\\
 \midrule
 \text{Tail beyond }u=2 & 2^{-144}\\
 \text{Midpoint quadrature for }f & 2^{-40}\\
 \text{Chebyshev source approximation} & 2^{-37}\\
 \text{Taylor heat approximation} & 2^{-116}\\
 \bottomrule
 \end{array}
\]
Their sum is below \(2^{-36}<10^{-8}\).

For \(I_\lambda\), comparison under the integral uses
\eqref{eq:delta} and the same bound on the cosine. Thus
\(\|H-I\|_{S,\boxD}\le S_*(T+2\e^{16}\delta)\).
This scalar baseline needs the same polynomial coefficients.
For \(z\ne0\), each exponential integral can, for example, be
defined by the finite identity
\[
 \int_0^2 P_\tau(u)\e^{\zeta u}\,\dd u
  =\left[\e^{\zeta u}\sum_{j=0}^{D_p}
                 \frac{(-1)^j P_\tau^{(j)}(u)}{\zeta^{j+1}}\right]_{u=0}^{2},
 \qquad \zeta=\pm\ii z .
\]
Differentiation telescopes because \(P_\tau^{(D_p+1)}=0\);
the half-sum of the two identities gives \(I_\lambda\).
\end{proof}

\section{Outstanding challenges}
\label{sec:challenges}

\paragraph{Coefficient acquisition and contraction.}
A computed version requires enclosures for the complete-source
coefficients, a stable choice of cores, and a bound for arithmetic
error in the oscillatory contraction. The theorem specifies exact
coefficients and exact arithmetic. Suppose initial node errors have
modulus at most \(e_{\rm node}\) in the unweighted polynomial values,
subsequent tensor compression has Euclidean error \(e_{\rm TT}\)
relative to that already perturbed weighted state, and the final
readout has additional scalar error at most \(e_{\rm out}\).
Their extra contribution to the scaled observation error is
bounded by
\[
 S_*\bigl(2\e^{16}e_{\rm node}
           +\sqrt{2}\,\e^{16}e_{\rm TT}+e_{\rm out}\bigr).
\]
The first transfer sums \(h\) times the uniform node error against
the cosine; the second uses Cauchy--Schwarz and
\(\sum_{\boldsymbol b}h|\cos(zu_{\boldsymbol b})|^2\le2\e^{32}\).
With \(\eta=10^{-8}/4\), sufficient budgets are respectively
\(\eta/2^{73}\), \(\eta/2^{73}\), and \(\eta/2^{48}\), because the
exact allowance is below \(\eta\).
For a differentiable approximate state \(\psi_{\rm a}\) evolved in
time by another method, the extra state error can be controlled
using \(r=\partial_\tau\psi_{\rm a}+K\psi_{\rm a}\).
Variation of constants gives the initial error propagated by
\(\e^{-\tau K}\) plus
\(\int_0^\tau\e^{-(\tau-s)K}r(s)\,\dd s\).
The propagator is contractive since its eigenvalues are
\(\e^{-\tau u_{\boldsymbol b}^2}\); hence the residual contribution
is at most \(\int_0^\tau\|r(s)\|_2\,\dd s\).
Demonstrating working precision and comparing total cost to the
scalar baseline require paying for acquisition and arithmetic in
both representations.

\paragraph{Growing heights.}
The constants in \eqref{eq:transfer} use a fixed upper height.
The Cauchy circle bound on the cosine increases with
\(|\operatorname{Re}z|\), and \(S(x)\) grows exponentially there.
Extending the argument requires a family of cutoffs, degrees
and bit depths with a quantitative error budget, including the
complete tail on the required scale. Even for the fixed box the
rank cap is a constructive upper bound. Determining smaller
ranks at specified precision, and their growth with the height,
remains a separate approximation problem.

\paragraph{Signed observations.}
The compact theorem controls the absolute error of the field.
A sign inference at a real point would require a lower bound on
the relevant scalar margin. Ratios and logarithms would additionally
require control near their zeros or denominators. Rank alone
cannot supply such a margin. For instance, the two-bit source with
coordinate \(u=b_1+b_2-1\) and amplitudes
\(\e^{\lambda-T_0},1,1,\e^{\lambda-T_0}\) gives the Fourier
contraction \(F_\lambda(z)=2+2\e^{\lambda-T_0}\cos z\).
It obeys \(\partial_\lambda F_\lambda=-\partial_z^2F_\lambda\).
At \(T_0=1/4\) and \(z=17\pi\) the state is a product with all
amplitudes one, while \(F_{T_0}=F_{T_0}'=0\) and \(F_{T_0}''=2\).
This atomic example has a positive source and a rank-one collision
at time and height coordinates inside \(\boxD\). The quantitative
bounds here use a different premise, the complete analytic theta
density in \eqref{eq:source}.

\paragraph{Acknowledgment.}
This draft was prepared with substantial OpenAI Codex LLM assistance.
The recorded drafting dispatch was GPT-6.1 Sol at Ultra reasoning effort.
Separate backend attestation is unavailable.

\begin{thebibliography}{9}
\bibitem{dlmf}
NIST Digital Library of Mathematical Functions,
\href{https://dlmf.nist.gov/5.11.E9}{Eq.~5.11.9},
\href{https://dlmf.nist.gov/20.2.E3}{Eq.~20.2.3}, and
\href{https://dlmf.nist.gov/20.7.E32}{Eq.~20.7.32}.

\bibitem{lindsey}
M.~Lindsey, \emph{Multiscale interpolative construction of quantized
tensor trains}, \href{https://arxiv.org/html/2311.12554v3}
{arXiv:2311.12554v3}, 2024.
\end{thebibliography}
\end{document}
